% GENERATED FILE. Edit canonical lessons, then run npm run generate:books.
% canonical-source-hash: fnv1a64:ab1f8369ddf37b4d
% canonical-lessons: PT-C500-costura, PT-C500-costurar, PT-C500-trico, PT-C500-tricotar, PT-C500-bordar

\chapter{Agulhas e linhas}
\label{ch:pt-agulhas-e-linhas}
\hlchaptermodality{500}

\begin{chapteropening}
\textbf{By the end of this chapter:} \emph{I can talk about sewing, knitting and embroidery.}

Complete the last lesson of chapter 500: I can talk about sewing, knitting and embroidery.
\end{chapteropening}

\section[a costura]{a costura — sewing; seam}

\label{lesson:PT-C500-costura}



\begin{quote}
Before the new one: say the Portuguese for faithful, then the Portuguese for faithfulness.
\end{quote}

\subsection*{You'll want to know: a costura}
\textbf{a costura} — \textquotedblleft{}sewing; seam\textquotedblright{}.

\begin{quote}
A minha avó aprendeu costura com a mãe.
\end{quote}

\begin{quote}
\emph{My grandmother learned sewing from her mother.}
\end{quote}

Also a seam: \emph{a costura rasgou}.

\begin{cousinweb}
\emph{costura} is usually traced to Vulgar Latin \emph{cōnsūtūra}, built on Latin \emph{cōnsuere} \textquotedblleft{}to sew together\textquotedblright{} (Lewis \& Short), from \emph{suere} \textquotedblleft{}to sew\textquotedblright{}.
\end{cousinweb}

\begin{grammarlens}[title={What you've built}]
One more word for telling, explaining and describing at B1.
\end{grammarlens}

\subsection*{Your turn}
\begin{itemize}
\raggedright
  \item \emph{Say it:} \emph{a costura}
  \item \emph{Say it:} \emph{a costura}, once more
  \item \emph{Say it:} a sentence of your own with \emph{a costura}
  \item \emph{Recall:} say the Portuguese for faithful, then the Portuguese for faithfulness, then say \emph{a costura} again
\end{itemize}

\begin{tcolorbox}[breakable,colback=teal!4,colframe=teal!35!black,title={Before you move on}]
What does a costura mean? (\textquotedblleft{}Sewing; seam\textquotedblright{}.) Say it once more.
\end{tcolorbox}
\section[costurar]{costurar — to sew}

\label{lesson:PT-C500-costurar}



\begin{quote}
Before the new one: say the Portuguese for faithfulness, then the Portuguese for sewing.
\end{quote}

\subsection*{You'll want to know: costurar}
\textbf{costurar} — \textquotedblleft{}to sew\textquotedblright{}.

\begin{quote}
Costurei um botão na camisa.
\end{quote}

\begin{quote}
\emph{I sewed a button on the shirt.}
\end{quote}

Built on \emph{a costura}. In Portugal \emph{coser} is also very common, especially for small mending jobs.

\begin{cousinweb}
\emph{costurar} is built on \emph{costura} \textquotedblleft{}sewing\textquotedblright{}, usually traced to Latin \emph{cōnsuere} \textquotedblleft{}to sew together\textquotedblright{} (Lewis \& Short).
\end{cousinweb}

\begin{grammarlens}[title={What you've built}]
One more word for telling, explaining and describing at B1.
\end{grammarlens}

\subsection*{Your turn}
\begin{itemize}
\raggedright
  \item \emph{Say it:} \emph{costurar}
  \item \emph{Say it:} \emph{costurar}, once more
  \item \emph{Say it:} a sentence of your own with \emph{costurar}
  \item \emph{Recall:} say the Portuguese for faithfulness, then the Portuguese for sewing, then say \emph{costurar} again
\end{itemize}

\begin{tcolorbox}[breakable,colback=teal!4,colframe=teal!35!black,title={Before you move on}]
What does costurar mean? (\textquotedblleft{}To sew\textquotedblright{}.) Say it once more.
\end{tcolorbox}
\section[o tricô]{o tricô — knitting}

\label{lesson:PT-C500-trico}



\begin{quote}
Before the new one: say the Portuguese for sewing, then the Portuguese for to sew.
\end{quote}

\subsection*{You'll want to know: o tricô}
\textbf{o tricô} — \textquotedblleft{}knitting\textquotedblright{}.

\begin{quote}
O tricô acalma-me depois do trabalho.
\end{quote}

\begin{quote}
\emph{Knitting calms me down after work.}
\end{quote}

\emph{Fazer tricô} — to knit. In Portugal \emph{fazer malha} is also common.

\begin{cousinweb}
\emph{tricô} was borrowed from French \emph{tricot} \textquotedblleft{}knitting\textquotedblright{} (EtymDB), from \emph{tricoter}.
\end{cousinweb}

\begin{grammarlens}[title={What you've built}]
One more word for telling, explaining and describing at B1.
\end{grammarlens}

\subsection*{Your turn}
\begin{itemize}
\raggedright
  \item \emph{Say it:} \emph{o tricô}
  \item \emph{Say it:} \emph{o tricô}, once more
  \item \emph{Say it:} a sentence of your own with \emph{o tricô}
  \item \emph{Recall:} say the Portuguese for sewing, then the Portuguese for to sew, then say \emph{o tricô} again
\end{itemize}

\begin{tcolorbox}[breakable,colback=teal!4,colframe=teal!35!black,title={Before you move on}]
What does o tricô mean? (\textquotedblleft{}Knitting\textquotedblright{}.) Say it once more.
\end{tcolorbox}
\section[tricotar]{tricotar — to knit}

\label{lesson:PT-C500-tricotar}



\begin{quote}
Before the new one: say the Portuguese for to sew, then the Portuguese for knitting.
\end{quote}

\subsection*{You'll want to know: tricotar}
\textbf{tricotar} — \textquotedblleft{}to knit\textquotedblright{}.

\begin{quote}
No inverno, a tia tricota camisolas para todos.
\end{quote}

\begin{quote}
\emph{In winter, our aunt knits jumpers for everyone.}
\end{quote}

Also \emph{fazer tricô}.

\begin{cousinweb}
\emph{tricotar} was borrowed from French \emph{tricoter} \textquotedblleft{}to knit\textquotedblright{} (EtymDB), which in Old French meant \textquotedblleft{}to beat with sticks\textquotedblright{}.
\end{cousinweb}

\begin{grammarlens}[title={What you've built}]
One more word for telling, explaining and describing at B1.
\end{grammarlens}

\subsection*{Your turn}
\begin{itemize}
\raggedright
  \item \emph{Say it:} \emph{tricotar}
  \item \emph{Say it:} \emph{tricotar}, once more
  \item \emph{Say it:} a sentence of your own with \emph{tricotar}
  \item \emph{Recall:} say the Portuguese for to sew, then the Portuguese for knitting, then say \emph{tricotar} again
\end{itemize}

\begin{tcolorbox}[breakable,colback=teal!4,colframe=teal!35!black,title={Before you move on}]
What does tricotar mean? (\textquotedblleft{}To knit\textquotedblright{}.) Say it once more.
\end{tcolorbox}
\section[bordar]{bordar — to embroider}

\label{lesson:PT-C500-bordar}



\begin{quote}
Before the new one: say the Portuguese for knitting, then the Portuguese for to knit.
\end{quote}

\subsection*{You'll want to know: bordar}
\textbf{bordar} — \textquotedblleft{}to embroider\textquotedblright{}.

\begin{quote}
A avó bordou flores num lenço.
\end{quote}

\begin{quote}
\emph{Grandma embroidered flowers on a handkerchief.}
\end{quote}

Embroidery is \emph{o bordado}.

\begin{cousinweb}
\emph{bordar} is linked in EtymDB with French \emph{broder} \textquotedblleft{}to embroider\textquotedblright{}, which goes back to a Germanic word for an edge or border.
\end{cousinweb}

\begin{grammarlens}[title={What you've built}]
One more word for telling, explaining and describing at B1.
\end{grammarlens}

\subsection*{Your turn}
\begin{itemize}
\raggedright
  \item \emph{Say it:} \emph{bordar}
  \item \emph{Say it:} \emph{bordar}, once more
  \item \emph{Say it:} a sentence of your own with \emph{bordar}
  \item \emph{Recall:} say the Portuguese for knitting, then the Portuguese for to knit, then say \emph{bordar} again
\end{itemize}

\begin{tcolorbox}[breakable,colback=teal!4,colframe=teal!35!black,title={Before you move on}]
What does bordar mean? (\textquotedblleft{}To embroider\textquotedblright{}.) Say it once more.
\end{tcolorbox}
